\chapter[Sermon 72. The Seventh Angel Pours Out His Vial.]{The Seventh Angel Pours Out His Vial.}
\label{sermon:072}
\markright{Sermon 72. The Seventh Angel Pours Out His Vial.}

And the seventh angel poured out his vial into the air. And there came a great voice out of Heaven from the throne, saying, “It is done.” And there followed voices, thunderings, and lightnings, and there was a great earthquake, such as was not since men were upon the earth, so mighty an earthquake and so great. And the great city was divided into three parts. And the cities of nations fell. And great Babylon came into remembrance before God, to give unto her the cup of wine of the fierceness of his wrath. And every island fled away, and mountains were not found. And there fell a great hail, as it had been talents, out of Heaven upon the men, and the men blasphemed God because of the hail, for it is great, and the plague of it was sore.

The seventh and last cup poured out into the air signifies the disturbance and alteration of all elements, and the horrible, but just judgment of God, and finally the end of all things and eternal punishment. The things are expressed with figurative language, taken mostly from the Prophets, and by a secret comparison brought from holy history. This is done for this consideration, that all things might be more majestic, and that every person should more diligently search for the meaning of an excellent matter, which once found, he might keep and retain in perfect memory.

And when the air is disturbed, diverse and horrible tempests arise in the air. And the Lord Jesus in the Gospel according to Matthew testifies that about the final coming of Christ, the powers of heaven shall be moved. And as soon as the cup was poured out into the air, and a great tempest arose, a voice sounded, it is done. By which voice is signified how all things are at an end, even of the whole world, much more of wicked papistry. And this voice is heard out of the very temple of heaven and the throne of God, lest we should doubt anything of the truth and certainty of the sentence given, and again of the virtue and power of him that pronounces it. Therefore are they shamefully deceived, so many as affirm the world to be everlasting, and that they shall reign always upon earth, and enjoy the pleasures thereof. A voice from heaven, out of the most holy temple of God, and even out of the most sincere throne of the Almighty, speaks, that it is done. For he speaks of the time to come as though it were past, that we might as certainly know that all worldly and popish things should have an end, as we undoubtedly know the things to be done which are already past. Let us therefore watch, and put no confidence in the things of this world, which are most deceptive. All things shall fall to decay and come to naught, men only, and the blessed spirits, remaining through the grace of God, the unhappy also remaining perpetually, appointed to perpetual punishment by the justice of God.

\index{John, Saint}And just as the holy Prophets depicted God’s judgment through figures for men’s eyes to see, so now the Lord Jesus, through Saint John, uses figurative language to shadow the terror of that horrible judgment. For he says that there will be thunderings, voices, lightning, and thunderbolts, and an earthquake so terrible that the world has never felt anything like it at any time. For Saint Peter also, at the end of his second letter, recounts terrible things about the last day and the burning of all worldly things. But the quaking and terror of men’s minds will be yet a great deal more terrible than all these.

The Lord in the Gospel according to Matthew. Then he says, “all the kindreds of the earth.” For the ungodly, whose consciences are wicked and corrupt, shall feel those terrors and torments unspeakable. The godly, as according to the saying of our Savior, do not come into judgment; so although that they also, by reason of the infirmity of the flesh, be somewhat astonished at the sudden alteration of things, and the terrible tearing and crashing of all elements, yet since they have known before that the same should come to pass, and believe the Savior’s saying, “your redemption draws near,” they gather up their spirits and comfort themselves in Christ, and rejoice in him, coming to judge or condemn the ungodly, but to save the godly. And herein is alluded to sundry stories of the holy scriptures, but chiefly to the burning of Sodom, to the drowning of Pharaoh in the Red Sea, and the ruin of Jericho, \&c. Those were verily but several destructions, and yet terrible above measure; therefore what think we that last destruction will be, which shall be general?

\index{Antichrist}\index{Papacy}Then shall that great city be cut asunder, the universality of men in the great church, divided into three parts: that is to say, in the end there shall be found three kinds of men in the church. There are true Christians, who attribute to Christ his true glory, that is, all things of true salvation, and cleave to him alone by sincere faith. There are Papists, who after the letter ascribe many things to Christ, but not as is fitting: for they ascribe to Antichrist those things which belong to Christ alone; and in associating with him, things that ought not to be associated, they deny Christ. For if the Pope is head of the universal church, if he is king and priest, and so forth, why is Christ preached to have those things alone? There are moreover Moderates, who will not seem to deny Christ, and yet attribute a little to Antichrist, whom nevertheless in many things they despise utterly. These have no certain religion, but one established and conceived at their pleasure, as it pleases them to believe this or that. There is a great number of these men at this day, deriding and mocking whatsoever is not tuned after their own light, and wanton Lucianicall wits. You may find also in the gospel a field sown with sundry seed, to bring forth most diverse fruits, yea even cockle and darnel, which at length in the end of the world shall be gathered. \&c. Matthew 13.

\index{Arethas of Caesarea}\index{Turks}Moreover, the cities of the Gentiles (he says) will fall, by which I understand the Jewish, Turkish, and foreign religions, torn into various sects or heresies. But each of these has its societies, rites, and laws, which they commend as the best, and as those that will endure forever; but they too will fall. The sole religion or faith of Christ shall prevail and overcome. Arethas expounding this passage in the same manner: The cities of the heathen, he says, falling down, are diverse opinions about faith, and so forth. They (I say) have all fallen.

\index{Rome}But especially, he affirms and demonstrates diligently, it was fitting and necessary that the city and church of Rome should be destroyed and subjected to perpetual torments. I have declared sufficiently before that Babylon is Rome, which in truth is great, not only in Italy but throughout all France, Spain, Germany, and other realms. The city and church of Rome have seemed to many to be everlasting and triumphant forever. Herein the Epicureans cry out that God does not care for these lesser things, but that every man lives here, either happily or unhappily, according as he has discreetly and skillfully framed his life, that [illegible] knows nothing of our pleasures and displeasures, and our conversation. But contrariwise, Saint John affirms that the Lord has remembered Babylon, and to have remembered her so that he has determined to commit her to torments. He expresses this through a prophetic phrase of speech, that he might give unto her the cup of wine of indignation, or fierceness of his wrath, that is to say, that he might punish her accordingly, as the great indignation and wrath of God requires. Therefore she shall have no small punishment, for the wrath of God is not light, but most grievous and hot. For he requites and recompenses the slackness of punishment with the extremity of pain and torment. You may read similar things in Malachi 3. Concerning the cup also, of the wine of God’s fury, has been spoken of before from the Prophets.

Now also, among other things, by a figurative speech he shows that the ungodly have no refuge, nor way to escape. Otherwise, the wealthier would hide themselves far away on islands, to be out of range of danger; many would flee into the mountains, that they might lurk safely. But now he says that even the islands flee, and therefore, in fleeing, they cannot be overtaken. He adds that the mountains—that is to say, no places of refuge or hiding—can be found. Therefore, nothing remains but that all the ungodly, in general, being taken should be put to torment.

\index{Primasius}Furthermore, he adds that hailstones as big as talents should be cast down from heaven upon wicked men, and that those who have not been remembered to have fallen in no memory of him [?]. And he seems to have alluded to the story of the Canaanites, which is in the tenth chapter of Joshua. To be short, here is signified that the grievous and inevitable judgment of God pronounced against all the ungodly shall, at the general judgment, torment the wicked with such an extremity that no eloquence of men, no sense nor understanding can attain unto, for it is always more grievous. Primasius, expounding this place, says: he sets the wrath of revenge in hail, whereof we read: “The wrath of the Lord falls down like hail.” Neither does he without cause mention a talent’s weight, for with equity will he inflict judgment, and so forth.

Here is shown the stubborn and incurable rebellion and impatience of the wicked, by which they are provoked against God’s judgments, and utter blasphemies against the Judge himself and his judgment. I have treated these things more briefly, for we have heard much the same before concerning the end of the 11th Chapter. To the Lord be praise and glory.
